\documentclass[12pt]{article}
\usepackage{amsmath, amssymb, amsthm}
\usepackage{geometry}
\usepackage{booktabs}
\usepackage{hyperref}
\geometry{margin=1.2in}

%% ---- theorem environments ----
\theoremstyle{definition}
\newtheorem{definition}{Definition}[section]
\theoremstyle{plain}
\newtheorem{theorem}{Theorem}[section]
\newtheorem{lemma}[theorem]{Lemma}
\theoremstyle{remark}
\newtheorem{remark}{Remark}[section]

%% ---- macros ----
\newcommand{\conf}[2]{\mathrm{conf}_{#1}(#2)}
\newcommand{\dN}[1]{\delta_1(N_{#1})}
\newcommand{\ddiff}{\delta_{\mathrm{diff}}^{*}}
\newcommand{\Xdom}{\mathcal{X}}
\newcommand{\Edom}{\mathcal{E}}

\title{%
  The Inequality $\delta_1(N_1)+\delta_{\mathrm{diff}}^{*} \geq \delta_1(N_2)$
  Does \textbf{Not} Hold in General:\\[6pt]
  \large A Counterexample with Full Calculations,
         and the Correct Inequality}
\author{}
\date{}

\begin{document}
\maketitle
\tableofcontents
\newpage

%% ====================================================================
\section{Setting}

Throughout this document:
\begin{itemize}
  \item Input domain: $\Xdom = [0,1]^2$, with $x=(x_1,x_2)$.
  \item Perturbation set: $\Edom = [-p,\,p]^2$ with $p=\tfrac{1}{10}$.
  \item $c_{\mathrm{tag}}=1$ (correct class), $c_{\mathrm{target}}=2$ (adversarial class).
  \item Attack margin $\eta = 10^{-3}$.
  \item Transition margin $\varepsilon_0 = 10^{-3}$ (used in constraint C1).
\end{itemize}

\begin{definition}[Confidence margin]\label{def:conf}
For a network $N$ with output logits
$\mathrm{logit}_N(x) = \bigl[\mathrm{logit}_N(x)[1],\,\mathrm{logit}_N(x)[2]\bigr]$,
\[
  \conf{N}{x}
  \;:=\;
  \mathrm{logit}_N(x)[1] \;-\; \mathrm{logit}_N(x)[2].
\]
$\conf{N}{x}>0$ means $N$ predicts class $c_{\mathrm{tag}}=1$ (correct prediction).
\end{definition}

\begin{definition}[$N$-foolable]\label{def:foolable}
$x\in\Xdom$ is \emph{$N$-foolable} if there exists $\varepsilon\in\Edom$ with
$x+\varepsilon\in\Xdom$ and
\[
  \mathrm{logit}_N(x+\varepsilon)[2]
  \;\geq\;
  \mathrm{logit}_N(x+\varepsilon)[1] + \eta.
\]
Equivalently, $\conf{N}{x+\varepsilon} \leq -\eta$.
\end{definition}

\begin{definition}[$\delta_1(N)$ — standard mode]\label{def:delta1}
\[
  \delta_1(N) := \sup\bigl\{\,\conf{N}{x}
    \;\big|\;
    x\in\Xdom,\;
    \conf{N}{x}\geq 0,\;
    x\ N\text{-foolable}\,\bigr\}.
\]
\end{definition}

\begin{definition}[$\ddiff$ — transfer MIP]\label{def:ddiff}
\[
  \ddiff := \sup\bigl\{\,
    \conf{N_2}{x}-\conf{N_1}{x}
    \;\big|\;
    \underbrace{\conf{N_1}{x}\geq\dN{1}+\varepsilon_0}_{(\mathrm{C1})},\;
    \underbrace{\conf{N_2}{x}\geq\conf{N_1}{x}}_{(\mathrm{C2})},\;
    \underbrace{x\ N_2\text{-foolable}}_{(\mathrm{C3})}
  \,\bigr\}.
\]
\end{definition}

%% ====================================================================
\section{Networks}

We use two \emph{linear} networks (the ReLU layers are inactive because
$x_1,x_2\in[0,1]\geq 0$):

\medskip
\begin{center}
\begin{tabular}{lll}
\toprule
Network & Logits & Confidence margin \\
\midrule
$N_1$ & $\mathrm{logit}_{N_1}(x)=\bigl[x_1,\;x_2\bigr]$
      & $\conf{N_1}{x} = x_1 - x_2$ \\[4pt]
$N_2$ & $\mathrm{logit}_{N_2}(x)=\bigl[2x_1,\;5x_2\bigr]$
      & $\conf{N_2}{x} = 2x_1 - 5x_2$ \\
\bottomrule
\end{tabular}
\end{center}

\medskip\noindent
Both are two-layer sequential nets with a $2\times2$ identity first layer,
an inactive ReLU, and a $2\times2$ diagonal second layer
($\mathrm{diag}(1,1)$ for $N_1$ and $\mathrm{diag}(2,5)$ for $N_2$).

%% ====================================================================
\section{Step-by-Step Calculations}

\subsection{Foolability conditions (closed form)}

\paragraph{$N_1$.}
$N_1(x+\varepsilon)$ predicts class 2 iff
$(x_2+\varepsilon_2)\geq(x_1+\varepsilon_1)+\eta$,
i.e., $\varepsilon_2-\varepsilon_1\geq(x_1-x_2)+\eta$.
The maximum of $\varepsilon_2-\varepsilon_1$ over $\Edom$ is
$p-(-p)=0.2$.
Hence:
\[
  \boxed{x\ N_1\text{-foolable}
  \;\Longleftrightarrow\;
  x_1-x_2 \;\leq\; 0.2-\eta.}
\]

\paragraph{$N_2$.}
$N_2(x+\varepsilon)$ predicts class 2 iff
$5(x_2+\varepsilon_2)\geq 2(x_1+\varepsilon_1)+\eta$,
i.e., $5\varepsilon_2-2\varepsilon_1\geq(2x_1-5x_2)+\eta$.
The maximum of $5\varepsilon_2-2\varepsilon_1$ over $\Edom$ is
$5\cdot 0.1+2\cdot 0.1=0.7$.
Hence:
\[
  \boxed{x\ N_2\text{-foolable}
  \;\Longleftrightarrow\;
  2x_1-5x_2 \;\leq\; 0.7-\eta.}
\]

\subsection{Computing $\dN{1}$}

Maximize $x_1-x_2$ subject to $x_1-x_2\leq 0.2-\eta$,
$x_1-x_2\geq 0$, $x\in[0,1]^2$:
\[
  \dN{1} = 0.2-\eta = 0.199.
\]
This is achieved at $x_1=0.199$, $x_2=0$:
\[
  \underbrace{x_1-x_2}_{=\,0.199}
  \leq 0.2-\eta=0.199. \quad\checkmark
\]

\subsection{Computing $\dN{2}$}

Maximize $2x_1-5x_2$ subject to $2x_1-5x_2\leq 0.7-\eta$,
$2x_1-5x_2\geq 0$, $x\in[0,1]^2$:
\[
  \dN{2} = 0.7-\eta = 0.699.
\]
Achieved at $x_2=0$, $x_1=\frac{0.699}{2}=0.3495$:
\[
  2(0.3495)-5(0) = 0.699 \leq 0.699.\quad\checkmark
\]

\noindent
\textbf{Explicit verification of $N_2$-foolability at $(0.3495,0)$
with $\varepsilon^*=(-0.1,+0.1)$:}
\[
  5(0+0.1)-2(0.3495-0.1)
  = 0.5 - 2\times 0.2495
  = 0.5 - 0.499
  = 0.001
  = \eta.
  \quad\checkmark
\]

\subsection{Computing $\ddiff$}

The transfer MIP objective is
\[
  \conf{N_2}{x}-\conf{N_1}{x}
  =(2x_1-5x_2)-(x_1-x_2)
  = x_1 - 4x_2.
\]
Constraints (substituting the closed forms from Section~3.1):
\begin{alignat}{2}
  &\text{C1:} & \quad x_1-x_2 &\;\geq\; \dN{1}+\varepsilon_0
    = 0.199+0.001 = 0.200, \label{eq:C1}\\
  &\text{C2:} & \quad x_1-4x_2 &\;\geq\; 0, \label{eq:C2}\\
  &\text{C3:} & \quad 2x_1-5x_2 &\;\leq\; 0.699, \label{eq:C3}\\
  &\text{domain:} & \quad x &\;\in\; [0,1]^2. \notag
\end{alignat}

\paragraph{Solving.}
Setting $x_2=0$ eliminates the $-4x_2$ penalty (the objective $x_1-4x_2$ is
maximised by setting $x_2$ as small as possible):
\begin{itemize}
  \item C1 becomes $x_1\geq 0.200$.
  \item C2 becomes $x_1\geq 0$ (automatically satisfied).
  \item C3 becomes $2x_1\leq 0.699$, i.e., $x_1\leq 0.3495$.
\end{itemize}
Maximising $x_1$ over $[0.200,\,0.3495]$ gives $x_1^*=0.3495$.

\[
  \ddiff = x_1^*-4\cdot 0 = \mathbf{0.3495}.
\]

\paragraph{Verification at $x^*=(0.3495,\,0)$:}
\[
\begin{array}{ll}
\text{C1:} & 0.3495-0=0.3495\geq 0.200.\quad\checkmark\\[2pt]
\text{C2:} & 0.3495-0=0.3495\geq 0.\quad\checkmark\\[2pt]
\text{C3:} & 2(0.3495)-5(0)=0.699\leq 0.699.\quad\checkmark\text{ (tight)}\\[4pt]
\multicolumn{2}{l}{\conf{N_1}{x^*}=0.3495,\quad
                   \conf{N_2}{x^*}=0.699,\quad
                   \text{difference}=0.3495=\ddiff.\quad\checkmark}
\end{array}
\]

%% ====================================================================
\section{The Stated Inequality Fails}

\begin{center}
\renewcommand{\arraystretch}{1.5}
\begin{tabular}{lc}
\toprule
Quantity & Value \\
\midrule
$\dN{1}$ (standard mode on $N_1$) & $0.199$ \\
$\ddiff$ (transfer MIP) & $0.3495$ \\
$\dN{1}+\ddiff$ & $\mathbf{0.5485}$ \\
\midrule
$\dN{2}$ (standard mode on $N_2$) & $\mathbf{0.699}$ \\
\bottomrule
\end{tabular}
\end{center}

\bigskip
\[
  \dN{1}+\ddiff
  \;=\; 0.5485
  \;\;\mathbf{<}\;\;
  0.699
  \;=\; \dN{2}.
\]

\noindent
The inequality $\dN{1}+\ddiff\geq\dN{2}$
\textbf{does not hold} for these networks. The gap is
$0.699-0.5485 = 0.1505$.

\paragraph{Why does the gap arise?}
The transfer MIP optimal point is $x^*=(0.3495,0)$ with
$\conf{N_1}{x^*}=0.3495$.
Since $\conf{N_1}{x^*}=0.3495 > 0.199=\dN{1}$, the transfer MIP's witness
has \emph{strictly higher} $N_1$-confidence than the worst-case $N_1$-vulnerable input.
Consequently:
\[
  \dN{2}
  \;\geq\; \conf{N_2}{x^*}
  \;=\; \underbrace{\conf{N_1}{x^*}}_{0.3495}
        + \underbrace{\ddiff}_{0.3495}
  \;=\; 0.699
  \;\;>\;\;
  \underbrace{\dN{1}}_{0.199}+\underbrace{\ddiff}_{0.3495}=0.5485.
\]
The formula $\dN{1}+\ddiff$ replaces $\conf{N_1}{x^*)=0.3495$ with the
\emph{smaller} quantity $\dN{1}=0.199$, losing the gap of
$0.3495-0.199=0.1505$.

%% ====================================================================
\section{The Correct Inequality}

The inequality holds in the \emph{reverse} direction.

\begin{theorem}\label{thm:correct}
If $\ddiff>-\infty$ (the transfer MIP is feasible), then
\[
  \boxed{\dN{2} \;\geq\; \dN{1} + \ddiff.}
\]
\end{theorem}

\begin{proof}
Fix $\epsilon>0$.
Let $x^*\in\Xdom$ be a point feasible for the transfer MIP achieving the
objective within $\epsilon$ of the supremum; write $\varepsilon^*$ for the
corresponding attack perturbation (from C3). Then:
\begin{align}
  \conf{N_1}{x^*} &\geq \dN{1}+\varepsilon_0, \tag{C1}\\
  \conf{N_2}{x^*} &\geq \conf{N_1}{x^*}, \tag{C2}\\
  x^* &\ N_2\text{-foolable (via }\varepsilon^*\text{)}, \tag{C3}\\
  \conf{N_2}{x^*}-\conf{N_1}{x^*} &\geq \ddiff-\epsilon. \tag{obj}
\end{align}

\noindent\textbf{Lower-bounding $\conf{N_2}{x^*}$:}
\begin{align*}
  \conf{N_2}{x^*}
  &= \conf{N_1}{x^*}
     + \bigl(\conf{N_2}{x^*}-\conf{N_1}{x^*}\bigr)\\
  &\geq (\dN{1}+\varepsilon_0)+(\ddiff-\epsilon)
  \tag{by C1 and (obj)}\\
  &\geq \dN{1}+\ddiff-\epsilon.
\end{align*}

\noindent\textbf{Feasibility for standard mode on $N_2$:}
By C2 and C1, $\conf{N_2}{x^*}\geq\conf{N_1}{x^*}\geq\dN{1}+\varepsilon_0>0$,
so $N_2$ \emph{correctly} classifies $x^*$.
Combined with C3 ($x^*$ is $N_2$-foolable), the point $x^*$ is feasible for
the standard-mode MIP on $N_2$. Therefore
\[
  \dN{2} \;\geq\; \conf{N_2}{x^*} \;\geq\; \dN{1}+\ddiff-\epsilon.
\]
Since $\epsilon>0$ was arbitrary, $\dN{2}\geq\dN{1}+\ddiff$.
\end{proof}

\paragraph{Numerical check (our networks):}
\[
  \dN{2}=0.699
  \;\geq\;
  \dN{1}+\ddiff=0.199+0.3495=0.5485.
  \quad\checkmark
\]
The gap of $0.1505=\conf{N_1}{x^*}-\dN{1}=0.3495-0.199$ is exactly the
surplus $N_1$-confidence at the transfer MIP optimum (cf.\ Remark~\ref{rem:gap}).

\begin{remark}[The gap quantified]\label{rem:gap}
Let $x^*$ denote the transfer MIP optimum. Then:
\[
  \dN{2}-\bigl(\dN{1}+\ddiff\bigr)
  \;\geq\;
  \conf{N_1}{x^*}-\dN{1}
  \;\geq\; \varepsilon_0 > 0.
\]
The gap is non-negative because C1 forces $\conf{N_1}{x^*}\geq\dN{1}+\varepsilon_0$,
so the transfer MIP's witness always has strictly higher $N_1$-confidence
than the worst-case $N_1$-vulnerable input.
\end{remark}

\end{document}
